% Figure (Paper B): how the analytical results depend on each other and where the algorithm uses them.
\begin{figure}[tbp]
\centering
\begin{tikzpicture}[
  node distance=5mm and 6mm,
  res/.style={draw=black!55, rounded corners=2pt, align=center, font=\footnotesize, inner sep=3pt, text width=26mm},
  key/.style={res, fill=blue!6, very thick},
  use/.style={font=\scriptsize\itshape, text=black!60, align=left},
  arr/.style={-{Stealth[length=2mm]}, black!60, thick}]
\node[res] (gcd) {\cref{lem:gcd}\\gcd splitting};
\node[res, right=of gcd] (ep) {\cref{lem:ep}\\finite dimension $\Leftrightarrow\EP$};
\node[res, right=of ep] (ident) {\cref{thm:ident}, \cref{prop:unique}\\identification};
\node[key, below=of ep] (factor) {\cref{thm:factor}\\groups are factors};
\node[res, left=of factor] (terms) {\cref{prop:terms}\\terms are not groups};
\node[res, right=of factor] (hm) {\cref{thm:hm}\\stability, $\gN$};
\node[res, below=of terms] (eint) {\cref{thm:eint}, \cref{prop:discont}\\interaction};
\node[res, below=of hm] (bayes) {\cref{prop:bayes-angle}\\angle sets error bars};
\node[res, below=of factor] (count) {\cref{prop:counting}\\counting condition};
\draw[arr] (ep)--(ident); \draw[arr] (gcd.north) -- ++(0,3mm) -| (ident.north);
\draw[arr] (gcd.south) -- ++(0,-1.5mm) -| (factor.north west);
\draw[arr] (ep)--(factor); \draw[arr] (ep)--(terms); \draw[arr] (factor)--(hm); \draw[arr] (terms)--(eint); \draw[arr] (hm)--(bayes);
\node[use, below=1mm of eint] {$\to$ $\rho$, reported only};
\node[use, below=1mm of count] {$\to$ label \FC};
\node[use, below=1mm of bayes] {$\to$ intervals, $P(\text{split})$};
\node[use, above=1mm of ident, xshift=4mm] {$\to$ no ``identified'' label};
\end{tikzpicture}
\caption{How the results of \cref{sec:theory} depend on each other (arrows) and where the algorithm uses them
(italic). The shaded result is the grouping principle; everything else either prepares it or measures how well it
can be carried out on finite data.}\label{fig:dependency}
\end{figure}
